\makeatletter
\ifx\documentclass\@twoclasseserror\else
\documentclass[danish,9pt,a5paper,twoside]{memoir}
\fi
\pagestyle{simple}
% Brug "empty" pagestyle til indholdsfortegnelsen:
% https://tex.stackexchange.com/a/140777/76220
\aliaspagestyle{index}{empty}
\aliaspagestyle{chapter}{empty}

\usepackage{float}
\usepackage{amsmath, amssymb, amsfonts}
\usepackage{verbatim}
\usepackage[utf8]{inputenc}
\usepackage[danish]{babel}
\usepackage{tabularx}
\usepackage{graphicx}
\usepackage{siunitx}
\usepackage{multicol}
\usepackage{relsize}
\usepackage{geometry}
\usepackage{wrapfig} %Used to make figures inline with text
\usepackage{hyperref}
\usepackage{microtype}
\usepackage[sc]{mathpazo} %sc=small caps, osf=sc+old style numerals
\usepackage[scaled=.95]{helvet}
\usepackage{courier}
\usepackage{svg}
\usepackage{adjustbox}
\usepackage{nicefrac}

% -------------------------------------
%  Dimensioner på siden
% -------------------------------------
\geometry{
    top=15mm,
    left=15mm,
    right=15mm,
    bottom=15mm
}

% -------------------------------------
% Kapitelopsætning
% -------------------------------------

\makechapterstyle{chapterstyle}{
    \renewcommand{\chapnamefont}{\normalfont\Large}
    \renewcommand{\chapnumfont}{\normalfont\Large}
    \renewcommand{\chaptitlefont}{\normalfont\Huge\bfseries}

    \renewcommand{\printchaptername}{}
    \renewcommand{\printchapternum}{}

    \renewcommand{\afterchapternum}{}
    \renewcommand{\afterchaptertitle}{\par\vspace{10pt}}
    \setlength{\beforechapskip}{0pt}
    \setlength{\afterchapskip}{15pt}
}

\chapterstyle{chapterstyle}

% -------------------------------------
% Afsnitsopsætning
% -------------------------------------

\setsecnumformat{}
\setsecheadstyle{\LARGE\bfseries}
\setsubsecheadstyle{\Large\bfseries}
\setsubsubsecheadstyle{\large\bfseries}

% -------------------------------------
% Sideopsætning
% -------------------------------------

\makepagestyle{pagestyle}

\makeoddhead{pagestyle}{}{\itshape\leftmark}{}
\makeevenhead{pagestyle}{}{\itshape\rightmark}{}

\makeoddfoot{pagestyle}{}{\thepage}{}
\makeevenfoot{pagestyle}{}{\thepage}{}

\pagestyle{pagestyle}
\aliaspagestyle{chapter}{pagestyle}
\renewcommand{\chaptermark}[1]{\markboth{#1}{#1}}

\setlength{\epigraphrule}{0pt}

\newcommand{\alkymia}{@lkymia }
\newcommand{\alkymias}{@lkymias }

% -------------------------------------
% Bestyrelses-tabeller
% -------------------------------------
\newlength{\boardrolewidth}
\setlength{\boardrolewidth}{1.1cm}

\newenvironment{board}{%
    \fontsize{8}{10}\selectfont
    \begin{tabular}{@{}p{\boardrolewidth}l@{}}
}{%
    \end{tabular}
    \par\bigskip
}
\newcommand{\boardtitle}[1]{%
    & \textbf{#1} \\
}
\newcommand{\bname}[1]{%
    \adjustbox{max width=\dimexpr\linewidth-\boardrolewidth\relax}{#1}%
}
\newenvironment{boardpage}{%
    \begin{multicols}{2}
}{%
    \end{multicols}
}

% -------------------------------------
% Citatbokse (svævende citater)
% -------------------------------------
% \citat[placering]{bredde}{citat}
% placering: r/l/i/o (som wrapfigure), default r
% bredde: andel af linewidth, fx 0.45\linewidth
\newcommand{\citat}[3][r]{%
    \begin{wrapfigure}{#1}{#2}
        \vspace{-4pt}
        \noindent\rule{\linewidth}{0.6pt}\par\vspace{4pt}
        \noindent\itshape\fontsize{8}{10}\selectfont``#3''\par\vspace{4pt}
        \noindent\normalfont\rule{\linewidth}{0.6pt}
        \vspace{-4pt}
    \end{wrapfigure}%
}
 
% \citatliste[placering]{bredde}{indhold}
% Samme som \citat, men til flere korte citater i forlængelse af hinanden.
% Adskil de enkelte citater med \\[4pt] i #3.
\newcommand{\citatliste}[3][r]{%
    \begin{wrapfigure}{#1}{#2}
        \vspace{-4pt}
        \noindent\rule{\linewidth}{0.6pt}\par\vspace{4pt}
        \noindent\itshape\fontsize{8}{10}\selectfont #3\par\vspace{4pt}
        \noindent\normalfont\rule{\linewidth}{0.6pt}
        \vspace{-4pt}
    \end{wrapfigure}%
}

% -------------------------------------
% Løbesedler og breve
% -------------------------------------
% \picture[bredde]{filnavn}
% Indsætter et billede centreret i en figure.
% "bredde" er andel af linewidth (eller en fast længde), default 0.55\linewidth.
\newcommand{\tkpicture}[2][0.55\linewidth]{%
    \begin{figure}[H]
        \centering
        \includegraphics[width=#1]{#2}
    \end{figure}%
}
